\section{Ethics, Generative AI Use, and Open Science}\label{sec:appendix}

\paragraph{Ethics considerations}
We analyze only public open-source GitHub repositories, statically: no server was executed and no live system or user was involved, and labeling was done by the authors. We report aggregate results without attributing intent to maintainers, and we describe individual servers by what they do rather than by name. For servers in which we observed behavior that may be security-relevant (Section~\ref{sec:cases}), notification of the maintainers is in progress at the time of submission; we will report its status in the final version of the paper.

\paragraph{Generative AI usage}
Generative AI is part of this work in two ways. (1) Object of study: we study the security of tool ecosystems used by LLM-based agents. (2) Research aid: an LLM-based assistant helped write analysis scripts, search literature, draft text, and summarize code changes for review; all labels were assigned by a human author. Generative AI was used for editorial purposes in this manuscript, and all outputs were inspected by the authors to ensure accuracy and originality. AI tools are not listed as authors.

\paragraph{Open science}
We release our code and data at \url{https://anonymous.4open.science/r/XXXX}\todo{anonymous artifact link}. The artifact contains the crawler, version sampler, tool and capability extractor, and the analysis and figure scripts; the list of the 2,571 servers with the commits we sampled; every detected change event; per-server summaries and every number reported in the paper; and the validation sample with its human labels. The per-version extraction snapshots (about 1.2\,GB) can be regenerated from the released code and will be archived on Zenodo upon acceptance.
